\section{2013秋}
\subsection{微分積分}

\prob{
  次の極限を求めよ。
  \begin{gather}
    \lim_{x \to \frac{1}{2}} \frac{\log 2x}{ 2x - 1 }
  \end{gather}
}

\prob{
  \begin{gather}
    \sin \bigl(x^2y^2\bigr) = x
  \end{gather}
  のとき、$\disp \frac{dy}{dx}$を求めよ。
}

\prob{
  次の定積分を求めよ。
  \begin{gather}
    \frac{1}{\pi} \int_{0}^{\pi} x\cos x \,dx
  \end{gather}
}


\prob{
  \begin{enumerate}[label=(\arabic*), ref=(\arabic*), itemsep=0pt]
    \item 以下の座標変換のヤコビアンを求めよ。
    \item 上記の座標変換を利用して、以下の積分の値を求めよ。
    ただし、$D \colon x^2+y^2+z^2 \leq 1$とする。
    \begin{gather}
      \iiint_{D} \sqrt{x^2 + y^2 + z^2} \,dxdydz
    \end{gather}
  \end{enumerate}
}

\prob{
  次の連立微分方程式を解け。
  \begin{gather}
    \begin{dcases*}
      \frac{d^2x}{dt^2} + 2x - y = 0 \\
      \frac{d^2y}{dt^2} - x + 2y = 0
    \end{dcases*}
  \end{gather}
}

\newpage
\subsection{線形代数}
\prob{
  未知数$x,\,y,\,z$についての次の連立方程式は、
  唯一の解を持たないが、$g = 3f - 2h$の場合には解が存在し得ることがわかっている。
  \begin{gather}
    % \begin{aligned}
      ax + bz = f \\
      3x + 2y + cz = g \\
      -y + dz = h
    % \end{aligned}
  \end{gather}
  \begin{enumerate}[label=(\arabic*), ref=(\arabic*), itemsep=0pt]
    \item $a$と$c$を$b,\,d,\,f,\,h$で表せ。
    \item 解を$b,\,d,\,f,\,h$と任意の定数を用いて求めよ。
  \end{enumerate}
}

\prob{
  \begin{enumerate}[label=(\arabic*), ref=(\arabic*), itemsep=0pt]
    \item\label{2013autumn:l2-1}
    次の2次形式
    \begin{gather}
      F(x,y) = 7x^2 + 8xy + y^2
    \end{gather}
    の標準形を求めよ。
    \item 次の式
    \begin{gather}
      G(x,y) = F(x,y) + 8x + 2y + 2 = 0
    \end{gather}
    は、2次曲線を表している可能性がある。
    問2.\ref{2013autumn:l2-1}の結果を用いて2次曲線になるかどうかを確認せよ。
    もし2次曲線である場合には、その形の名称を示した上で、問2.\ref{2013autumn:l2-1}の2次形式の主軸を用いた略図を示せ。
  \end{enumerate}
}

